\documentclass[11pt]{article}

% 基础包
\usepackage[utf8]{inputenc}
\usepackage[T1]{fontenc}
\usepackage{amsmath, amssymb, amsfonts}
\usepackage{graphicx}
\usepackage{hyperref}
\usepackage{natbib}
\usepackage{geometry}
\geometry{margin=1in}

% 标题
\title{Communication-Efficient Federated Learning with Subspace Variance Reduction}
\author{AutoSurvey Generated}
\date{\today}

\begin{document}

\maketitle

\section{Related Work}

\subsection{Foundational Algorithms and the Challenge of Heterogeneity in Federated Learning}
The seminal Federated Averaging (FedAvg) algorithm established the core communication-computation trade-off in federated learning (FL) by amortizing communication costs via multiple local stochastic gradient descent (SGD) steps \citep{konevcny2016federated}. Early analyses proved FedAvg can achieve convergence under non-IID data \citep{li2019convergence} and even linear speedup with partial participation in non-convex settings \citep{Yang2021AchievingLS}. However, the core pathology of \textit{client drift}—where local updates diverge from the global objective due to heterogeneous data—severely degrades performance and stability \citep{li2019convergence, su2023non}. This drift motivates the development of robust correction mechanisms. The SCAFFOLD algorithm introduced control variates to correct biased local updates, providing convergence guarantees robust to arbitrary heterogeneity and partial participation \citep{karimireddy2020scaffold}. Enhancements integrating momentum further improve convergence rates for both FedAvg and SCAFFOLD under unbounded heterogeneity \citep{Cheng2023MomentumBN}. Other variants address system imperfections: Federated Postponed Broadcast (FedPBC) mitigates bias from non-uniform client participation \citep{xiang2023towards}, and CC-FedAvg accommodates heterogeneous computational resources \citep{Zhang2022CCFedAvgCC}. Architectures like Ψ-Net enforce structural alignment across clients to handle heterogeneity \citep{Yu2020HeterogeneousFL}. Despite these advances, a fundamental dilemma persists: increasing local computation reduces communication rounds but inherently amplifies client drift, while robust correction methods like SCAFFOLD maintain or increase per-round communication bandwidth \citep{karimireddy2020scaffold}. This unresolved tension between local efficiency and heterogeneity-induced performance loss underscores the need for techniques that simultaneously reduce both communication frequency and per-round payload size without sacrificing robustness.

\subsection{Direct Communication Compression Techniques and Their Limits Under Heterogeneity}
A primary approach to reduce per-round bandwidth employs direct gradient compression. Techniques are categorized by their statistical properties: unbiased compressors (e.g., random sparsification, stochastic quantization) preserve the expected gradient direction but inject additional variance \citep{konevcny2016federated, He2023UnbiasedCS}, while biased or contractive compressors (e.g., Top-K sparsification, deterministic quantization) offer higher compression ratios but introduce systematic directional error \citep{Haddadpour2020FederatedLW}. The Error Feedback (EF) framework is pivotal for correcting compression bias by accumulating residuals in memory, stabilizing convergence with aggressive compressors \citep{Tang2019DoubleSqueezePS, Fatkhullin2023MomentumPI, Qian2020ErrorCD}. Unified analyses provide sharp convergence guarantees for FL algorithms combining local SGD with compression \citep{Haddadpour2020FederatedLW, basu2019qsparse}. Frameworks like LoCoDL integrate local training with compression for doubly-accelerated efficiency \citep{Condat2024LoCoDLCD}, and SCALLION extends SCAFFOLD's variance reduction to support unbiased compression \citep{Huang2023StochasticCA}. However, these methods face inherent limits under heterogeneity. Compression distortion interacts destructively with client drift, slowing convergence or requiring stringent assumptions \citep{li2019convergence, basu2019qsparse}. Algorithm-agnostic lower bounds reveal intrinsic performance costs tied to the compressor's parameter \citep{He2023LowerBA}. Furthermore, unbiased compression alone may not reduce total communication cost without independent compressors across workers \citep{He2023UnbiasedCS}. Decentralized settings show convergence degradation with poor network connectivity and aggressive compression \citep{Koloskova2019DecentralizedSO}. While stochastic distributed learning methods support general compression with variance reduction \citep{horvath2023stochastic}, their analyses often rely on convex or strong convexity assumptions \citep{horvath2023stochastic, zhou2021communication}. High-performance sparse communication layers like SparCML facilitate deployment \citep{Renggli2018SparCMLHS}, but the core limitation remains: direct compression reduces payload size but can amplify noise and bias under non-IID data, failing to holistically address the dual challenge of efficiency and robustness.

\subsection{Subspace and Low-Dimensional Learning Methods for Intrinsic Efficiency}
Subspace and low-rank methods pursue intrinsic efficiency by constraining optimization to a low-dimensional manifold, drastically reducing the number of parameters communicated per round. FedSub employs primal-dual subspace projection to reduce communication, computation, and memory costs while providing non-convex convergence guarantees \citep{zhang2025efficient}. Federated low-rank adaptation (LoRA) frameworks, such as ILoRA, learn low-rank updates to a base model, addressing challenges like misaligned initialization, rank-heterogeneous aggregation, and exacerbated client drift under non-IID data \citep{Zhou2025ILoRAFL}. FedMuon applies matrix orthogonalization to local updates for structure-aware optimization \citep{liu2025fedmuon}. FedLoRU uses accumulated low-rank client updates, connecting to implicit regularization benefits \citep{Park2024CommunicationEfficientFL}. Theoretical analyses show that FedAvg with local updates can learn a shared linear representation across clients \citep{collins2022fedavg}, and that widening overparameterized networks mitigates heterogeneity impact in the Neural Tangent Kernel regime \citep{jian2025widening}. However, standalone subspace methods exhibit critical gaps under heterogeneity. Operating in a strict subspace introduces a residual bias if the optimal direction lies outside the constrained space \citep{zhang2025efficient}. Random subspace initialization causes instability, and rank incompatibility during aggregation introduces errors \citep{Zhou2025ILoRAFL}. The subspace constraint can exacerbate client drift by preventing clients from fully correcting their local gradients toward the global objective \citep{Zhou2025ILoRAFL, zhang2025efficient}. Convergence analyses often rely on bounded gradient assumptions or idealized conditions like population gradients and Gaussian features \citep{zhang2025efficient, collins2022fedavg}. Comparative evaluations under controlled heterogeneity highlight performance variability \citep{reguieg2023comparative}. These limitations demonstrate that while subspace methods achieve intrinsic communication compression, they lack explicit mechanisms to correct the bias and variance induced by data heterogeneity, leading to unstable convergence. This gap motivates integrating subspace constraints with principled variance reduction to simultaneously achieve low per-round bandwidth and robustness to non-IID data.

\bibliographystyle{plainnat}
\bibliography{references}

\end{document}
